\ea\label{ex:text23-1}
A varigomo eove o peho kaku vagana vaa Teapu to paku raara bona kave.\\
\gll a	varigomo	eove	o	peho	kaku	vagana	vaa	teapu	to	paku	ra=ara	bona	kave\\
     A	varigomo	\textsc{h}.3\textsc{sg}.\textsc{pron}	\textsc{art}3.\textsc{sg}	\textsc{indef}3.\textsc{sg}	kind	fish	\textsc{lnk}	Teapu	\textsc{rel}	do	1\textsc{pl}.\textsc{in}.\textsc{ipfv}=1\textsc{pl}.\textsc{in}.\textsc{pron}	\textsc{obj}.\textsc{art}2.\textsc{sg}	net\\
\glt `Varigomo is a kind of fishing in the Teop area that we do with a net.'
\z

\ea\label{ex:text23-2}
A kave nae he a bebeahu ae a mata rutaa.\\
\gll a	kave	nae	he	a	bebeahu	ae	a	mata	rutaa\\
     \textsc{art}2.\textsc{sg}	net	3\textsc{sg}.\textsc{poss}.\textsc{pron}	\textsc{advers}	\textsc{art}2.\textsc{sg}	long	and1	\textsc{art}2.\textsc{sg}	mesh	small\\
\glt `The net used for it is long and has small meshes.'
\z

\ea\label{ex:text23-3}
O kaku vagana bona na pakupaku raara teo gohoho ge teo tarahi, ei tea maa hana taataana sa ei haana namana.\\
\gll o	kaku	vagana	bona	na	paku{$\sim$}paku	ra=ara	te=o	gohoho	ge	te=o	tarahi	ei	te=a	maa	hana	taa{$\sim$}taana	sa	ei	haa=na	namana\\
     \textsc{art}3.\textsc{sg}	kind	fish	\textsc{ana}.\textsc{sg}	\textsc{tm}2	\textsc{red}{$\sim$}do	1\textsc{pl}.\textsc{in}.\textsc{ipfv}=1\textsc{pl}.\textsc{in}.\textsc{pron}	\textsc{prep}=\textsc{art}3.\textsc{sg}	incoming.tide	or	\textsc{prep}=\textsc{art}3.\textsc{sg}	ebbing.tide	\textsc{dem}2	\textsc{prep}=\textsc{art}2.\textsc{sg}	\textsc{plm}	path	\textsc{red}{$\sim$}walk.on.the.reef	\textsc{neg}1	\textsc{dem}2	\textsc{neg}2=3\textsc{sg}.\textsc{ipfv}	ocean\\
\glt `We do this kind of fishing during the incoming tide or the ebbing tide, this means at places where we can walk, this is not in the deep ocean.'\footnote{lit. `This kind of fishing we do ...'}
\z

\ea\label{ex:text23-4}
O kaku vagana bona sa vaapenapena haana, beo aba o bero nao, evehee a bero amaa vuri, o kukan koa o sinivi to tahoa rori.\\
\gll o	kaku	vagana	bona	sa	vaa-penapena	haa=na	be=o	aba	o	bero	nao	evehee	a	bero	a=maa	vuri	o	kukan	koa	o	sinivi	to	ta-hoa	rori\\
     \textsc{art}3.\textsc{sg}	kind	fish	\textsc{ana}.\textsc{sg}	\textsc{neg}1	\textsc{caus}-be.forbidden	\textsc{neg}2=3\textsc{sg}.\textsc{ipfv}	\textsc{compl}2=\textsc{art}2.\textsc{pl}	person	\textsc{art}2.\textsc{pl}	many	go	but	\textsc{art}2.\textsc{sg}	many	\textsc{art}2.\textsc{sg}=\textsc{plm}	time	\textsc{art}3.\textsc{sg}	three	only	\textsc{art}3.\textsc{sg}	canoe	\textsc{rel}	\textsc{acaus}-push	3\textsc{pl}.\textsc{ipfv}\\
\glt `This kind of fishing does not forbid that many people go, but most of the time it is only three canoes that set off.'
\z

\ea\label{ex:text23-5}
O peho sinivi to vaahae nana bona kave mi bona kukan a aba, a buaku vaa tea vavaagunaha kave, a peha he a too horohoro.\\
\gll o	peho	sinivi	to	vaa-hae	nana	bona	kave	mi	bona	kukan	a	aba	a	buaku	vaa	tea	va{$\sim$}vaa-gunaha	kave	a	peha	he	a	too	horo{$\sim$}horo\\
     \textsc{art}3.\textsc{sg}	one3	canoe	\textsc{rel}	\textsc{caus}-board	3\textsc{sg}.\textsc{ipfv}	\textsc{obj}.\textsc{art}2.\textsc{sg}	net	with	\textsc{obj}.\textsc{art}2.\textsc{sg}	three	\textsc{art}2.\textsc{sg}	person	\textsc{art}2.\textsc{sg}	two	\textsc{lnk}	\textsc{compl}1	\textsc{red}{$\sim$}\textsc{caus}-go.down	net	\textsc{art}2.\textsc{sg}	one2.\textsc{sg}	\textsc{advers}	\textsc{art}2.\textsc{sg}	person	\textsc{red}{$\sim$}punt\\
\glt `One canoe takes the net with three people on board, two for letting the net down, and one person is the punter.'\footnote{lit. `One canoe that takes the net with three people on board, ...'}
\z

\ea\label{ex:text23-6}
O meho buo sinivi he mene vahakaa.\\
\gll o	meho	buo	sinivi	he	mene	vahakaa\\
     \textsc{art}3.\textsc{sg}	other3.\textsc{sg}	two3.\textsc{sg}	canoe	\textsc{advers}	for	beat.water\\
\glt `But the other two canoes have to beat the water (to chase the fish into the net).'\footnote{lit. `The other two canoes, however, (are) for beating (the) water.'}
\z

\ea\label{ex:text23-7}
O kaku vagana nae he na paku vonaen:\\
\gll o	kaku	vagana	nae	he	na	paku	vo=na=en\\
     \textsc{art}3.\textsc{sg}	kind	fish	3\textsc{sg}.\textsc{poss}.\textsc{pron}	\textsc{advers}	\textsc{tm}2	do	like=3\textsc{sg}.\textsc{ipfv}=\textsc{dem}3\\
\glt `The fishing with it is done as follows:'
\z

\ea\label{ex:text23-8}
Beara hua batara, eara toro tara bata naora beo iana an rori.\\
\gll be=ara	hua	bata=ra	e-ara	toro	tara	bata	nao=ra	be=o	iana	an	rori\\
     when=1\textsc{pl}.\textsc{in}.\textsc{pron}	paddle	\textsc{cont}=1\textsc{pl}.\textsc{in}.\textsc{ipfv}	\textsc{h}-1\textsc{pl}.\textsc{in}.\textsc{pron}	\textsc{tm}7	see	\textsc{cont}	thither=1\textsc{pl}.\textsc{in}.\textsc{ipfv}	\textsc{compl}2=\textsc{art}2.\textsc{pl}	fish	eat	3\textsc{pl}.\textsc{ipfv}\\
\glt `While we are paddling along, we must look where the fish are eating.'
\z

\ea\label{ex:text23-9}
Bea too horohoro sue voen, na mataa nana, eara repaa hua vahuus o pata iana bari.\\
\gll be=a	too	horo{$\sim$}horo	sue	vo=en	na	mataa	nana	e-ara	re=paa	hua	vahuus	o	pata	iana	bari\\
     when=\textsc{art}2.\textsc{sg}	person	\textsc{red}{$\sim$}punt	say	like=\textsc{dem}3	\textsc{tm}2	good	3\textsc{sg}.\textsc{ipfv}	\textsc{h}-1\textsc{pl}.\textsc{in}.\textsc{pron}	\textsc{consec}=\textsc{tm}3	paddle	come.close.to	\textsc{art}3.\textsc{sg}	school	fish	\textsc{ana}.\textsc{pl}\\
\glt `When the punter says it is good, we paddle close to that school of fish.'
\z

\ea\label{ex:text23-10}
A peha aba repaa gunaha me bono peho punete kave.\\
\gll a	peha	aba	re=paa	gunaha	me	bono	peho	punete	kave\\
     \textsc{art}2.\textsc{sg}	one2.\textsc{sg}	person	\textsc{consec}=\textsc{tm}3	go.down	with	\textsc{obj}.\textsc{art}3.\textsc{sg}	one3.\textsc{sg}	edge	net\\
\glt `Then one person goes down with one edge of the net.'
\z

\ea\label{ex:text23-11}
O sinivi repaa hua taviri bono pata iana, bea meha aba rurun gunaha bata ni nana bona kave.\\
\gll o	sinivi	re=paa	hua	taviri	bono	pata	iana	be=a	meha	aba	rurun	gunaha	bata	ni	nana	bona	kave\\
     \textsc{art}3.\textsc{sg}	canoe	\textsc{consec}=\textsc{tm}3	paddle	go.around	\textsc{obj}.\textsc{art}3.\textsc{sg}	school	fish	when=\textsc{art}2.\textsc{sg}	other2.\textsc{sg}	person	drop	go.down	\textsc{cont}	\textsc{appl}	3\textsc{sg}.\textsc{ipfv}	\textsc{obj}.\textsc{art}2.\textsc{sg}	net\\
\glt `The canoe then paddles around the school of fish, while the other man is dropping the net along the way.'
\z

\ea\label{ex:text23-12}
O meho buo sinivi he paa taneo vaevuru tea vahakaa, o meho tea papa matau, o meho he tea papa keruka.\\
\gll o	meho	buo	sinivi	he	paa	taneo	vaevuru	tea	vahakaa	o	meho	te=a	papa	matau	o	meho	he	te=a	papa	keruka\\
     \textsc{art}3.\textsc{sg}	other3.\textsc{sg}	two3.\textsc{sg}	canoe	\textsc{advers}	\textsc{tm}3	start	already	\textsc{compl}1	beat.water	\textsc{art}3.\textsc{sg}	other3.\textsc{sg}	\textsc{prep}=\textsc{art}2.\textsc{sg}	side	right	\textsc{art}3.\textsc{sg}	other3.\textsc{sg}	\textsc{advers}	\textsc{prep}=\textsc{art}2.\textsc{sg}	side	left\\
\glt `The other two canoes have already started to beat the water, one on the right side, the other one on the left side.'
\z

\ea\label{ex:text23-13}
Pasi vatete maari bono pata iana, ore vatete hovo komana kave.\\
\gll pasi	vatete	maa=ri	bono	pata	iana	o=re	vatete	hovo	koma-n=a	kave\\
     \textsc{tm}4	chase	hither=3\textsc{pl}.\textsc{ipfv}	\textsc{obj}.\textsc{art}3.\textsc{sg}	school	fish	3\textsc{pl}.\textsc{pron}=\textsc{consec}	chase	enter	inside-3\textsc{sg}.\textsc{poss}=\textsc{art}2.\textsc{sg}	net\\
\glt `They will chase the school of fish, and then chase them into the net,'
\z

\ea\label{ex:text23-14}
Bea bua too vavaagunaha kave rahi vatatana bata ni rori bono buo punete kave, eori repaa asun bona iana en komana kave.\\
\gll be=a	bua	too	va{$\sim$}vaa-gunaha	kave	rahi	vatatana	bata	ni	rori	bono	buo	punete	kave	e-ori	re=paa	asun	bona	iana	en	koma-n=a	kave\\
     when=\textsc{art}2.\textsc{sg}	two2.\textsc{sg}	person	\textsc{red}{$\sim$}\textsc{caus}-go.down	net	pull	meet	\textsc{cont}	\textsc{appl}	3\textsc{pl}.\textsc{ipfv}	\textsc{obj}.\textsc{art}3.\textsc{sg}	two3.\textsc{sg}	edge	net	\textsc{h}-3\textsc{pl}.\textsc{pron}	\textsc{consec}=\textsc{tm}3	kill	\textsc{obj}.\textsc{art}2.\textsc{sg}	fish	\textsc{dem}3	inside-3\textsc{sg}.\textsc{poss}=\textsc{art}2.\textsc{sg}	net\\
\glt `When the two net-droppers pull the two edges of the net together, they kill the fish inside the net.'
\z

\ea\label{ex:text23-15}
Amaa aba vai to sa paa hana hau tea maa mata kave, kahi tahitahi ri vuara bona makaree.\\
\gll a=maa	aba	vai	to	sa	paa	hana	ha=u	te=a	maa	mata	kave	kahi	tahi{$\sim$}tahi	ri	vu=ara	bona	makaree\\
     \textsc{art}2.\textsc{sg}=\textsc{plm}	one	\textsc{dem}6	\textsc{rel}	\textsc{neg}1	\textsc{tm}3	hang	\textsc{neg}2=\textsc{pfv}	\textsc{prep}=\textsc{art}2.\textsc{sg}	\textsc{plm}	mesh	net	\textsc{tm}1	\textsc{red}{$\sim$}spear	3\textsc{pl}.\textsc{objm}	\textsc{pfv}=1\textsc{pl}.\textsc{in}.\textsc{pron}	\textsc{obj}.\textsc{art}2.\textsc{sg}	fishing.spear\\
\glt `Those that have not been caught in the meshes of the net, we would spear with the fishing spear.'
\z

\ea\label{ex:text23-16}
Amaa aba vai to paa hanau teo mata kave na asun ri raara bono popoura.\\
\gll a=maa	aba	vai	to	paa	hana=u	te=o	mata	kave	na	asun	ri	ra=ara	bono	popoura\\
     \textsc{art}2.\textsc{sg}=\textsc{plm}	one	\textsc{dem}6	\textsc{rel}	\textsc{tm}3	hang=\textsc{pfv}	\textsc{prep}=\textsc{art}2.\textsc{pl}	mesh	net	\textsc{tm}2	kill	3\textsc{pl}.\textsc{objm}	1\textsc{pl}.\textsc{in}.\textsc{ipfv}=1\textsc{pl}.\textsc{in}.\textsc{pron}	\textsc{obj}.\textsc{art}3.\textsc{sg}	baton\\
\glt `Those that have been caught in the meshes of the net, we kill them with the baton.'
\z

\newpage
\ea\label{ex:text23-17}
Be kavara, eara repaa vaahae vahabana a kave, eara repaa nao pete teo meho hum vai to anan bata rio iana.\\
\gll be	kavara	e-ara	re=paa	vaa-hae	vahabana	a	kave	e-ara	re=paa	nao	pete	te=o	meho	hum	vai	to	an{$\sim$}an	bata	ri=o	iana\\
     when	finished	\textsc{h}-1\textsc{pl}.\textsc{in}.\textsc{pron}	\textsc{consec}=\textsc{tm}3	\textsc{caus}-board	again	\textsc{art}2.\textsc{sg}	net	\textsc{h}-1\textsc{pl}.\textsc{in}.\textsc{pron}	\textsc{consec}=\textsc{tm}3	go	subsequently	\textsc{prep}=\textsc{art}3.\textsc{sg}	other3.\textsc{sg}	place	\textsc{dem}6	\textsc{rel}	\textsc{red}{$\sim$}eat	\textsc{cont}	3\textsc{pl}.\textsc{ipfv}=\textsc{art}2.\textsc{pl}	fish\\
\glt `After that we take the net again on board and go on to another place where the fish eat.'\footnote{lit. `When finished, then we take the net again on board, then we subsequently go to another place ...'}
\z

\ea\label{ex:text23-18}
A vaakavara vagana na tei bata koana tea masi ge beo gohoho beera oha maa, eara repaa voosu vaha vaan.\\
\gll a	vaa-kavara	vagana	na	tei	bata	koa=na	te=a	masi	ge	be=o	gohoho	beera	oha	maa	e-ara	re=paa	voosu	vaha	vaan\\
     \textsc{art}2.\textsc{sg}	\textsc{caus}-finished	fish	\textsc{tm}2	be.there	\textsc{cont}	just=3\textsc{sg}.\textsc{ipfv}	\textsc{prep}=\textsc{art}2.\textsc{sg}	low.tide	or	when=\textsc{art}3.\textsc{sg}	incoming.tide	big	exceed	hither	\textsc{h}-1\textsc{pl}.\textsc{in}.\textsc{pron}	\textsc{consec}=\textsc{tm}3	go.home	back	village\\
\glt `The end of the fishing trip is at low tide, and when the incoming tide is too big, then we go home, back to the village.'
\z
